\chapter{Absolute values, valuations, and graded reduction}
\index{valuation}\index{graded reduction}
\label{chap:valuations}

Valuations occur in this book in three roles. They are points of analytic
spaces, coordinates on skeleta, and probes of singularities. These roles use
slightly different languages. This chapter fixes the dictionary and, in
particular, the sign convention that prevents later formulae from changing
direction unexpectedly.

The chapter proceeds from the elementary to the structural. We first record
the equivalence between multiplicative seminorms and real semivaluations,
then introduce monomial and quasimonomial valuations, completed residue
fields, and graded reduction. Valuation rings, centers, and the extension
theory of valuations along field extensions follow. The numerical invariants
of a valuation---rank, rational rank, and residue transcendence
degree---are then tied together by Abhyankar's inequality, and a gallery of
fully computed valuations on $k(x,y)$ realizes every phenomenon the
inequality permits. The chapter closes with the linear-growth condition used
in \cref{part:bj}, a map from each class of valuations to its later role,
and the multiplicative--additive dictionary.

Throughout, rings are commutative with unit. Unless an ordered group is
explicitly mentioned, a valuation is real-valued. We put
$\overline{\mathbf R}=\mathbf R\cup\{+\infty\}$ with
$a+(+\infty)=+\infty$.

\section{Absolute values and semivaluations}

\begin{definition}\label{def:absolute-value}
An \emph{absolute value} on a field $K$ is a map
$|\cdot|:K\to\mathbf R_{\geq0}$ satisfying
\[
 |a|=0\Longleftrightarrow a=0,\quad |ab|=|a||b|,\quad
 |a+b|\leq |a|+|b|.
\]
It is \emph{non-Archimedean} if
\[
                         |a+b|\leq\max\{|a|,|b|\}.
\]
The absolute value is \emph{trivial} if $|a|=1$ for every $a\ne0$.
\end{definition}

\begin{definition}\label{def:multiplicative-seminorm-semivaluation}
Let $A$ be a ring. A \emph{multiplicative seminorm} on $A$ is a map
$|\cdot|:A\to\mathbf R_{\geq0}$ with $|0|=0$, $|1|=1$,
$|fg|=|f||g|$, and $|f+g|\leq\max\{|f|,|g|\}$. A
\emph{real semivaluation} on $A$ is a map
\[
                    v:A\longrightarrow\overline{\mathbf R}
\]
such that
\[
 v(0)=+\infty,\quad v(1)=0,\quad
 v(fg)=v(f)+v(g),\quad
 v(f+g)\geq\min\{v(f),v(g)\}.
\]
The two notions correspond by
\[
                         v=-\log|\cdot|,
             \qquad |f|=\rme^{-v(f)},
\]
where $\rme^{-\infty}=0$.
\end{definition}

\begin{lemma}\label{lem:seminorm-semivaluation-dictionary}
\index{seminorm!multiplicative}\index{semivaluation}
The assignments $v=-\log|\cdot|$ and $|\cdot|=\rme^{-v}$ are mutually
inverse bijections between multiplicative seminorms and real semivaluations
on a ring $A$. Under this bijection,
\[
 |f|\leq1\Longleftrightarrow v(f)\geq0,
 \qquad
 |f|<1\Longleftrightarrow v(f)>0,
 \qquad
 |f|=0\Longleftrightarrow v(f)=+\infty.
\]
For $s>0$, the seminorm $|\cdot|^s$ corresponds to the semivaluation $sv$.
\end{lemma}
\begin{proof}
The function $-\log:[0,+\infty)\to\overline{\mathbf R}$, with
$-\log0=+\infty$, is a strictly decreasing bijection converting products
into sums and $\max$ into $\min$. Each axiom of a multiplicative seminorm is
carried to the corresponding axiom of a semivaluation and conversely:
$|0|=0$ becomes $v(0)=+\infty$; $|1|=1$ becomes $v(1)=0$;
$|fg|=|f||g|$ becomes $v(fg)=v(f)+v(g)$; and
$|f+g|\leq\max\{|f|,|g|\}$ becomes $v(f+g)\geq\min\{v(f),v(g)\}$. The three
displayed equivalences and the scaling statement are immediate from
monotonicity and from $-\log(t^s)=-s\log t$.
\end{proof}

The base $\rme$ of the exponential is a normalization, not a structure: any
base $>1$ produces the same theory after rescaling $v$. This innocuous
observation returns with force over a trivially valued field, where no
element of the ground field pins the scale down and rescaling becomes the
positive flow on the analytification; see
\cref{chap:trivial-analytification}.

The next lemma separates the two regimes in the first definition. It is
the reason that non-Archimedean behaviour can be detected on the prime
field alone.

\begin{lemma}\label{lem:ultrametric-integer-bound}
\index{absolute value!non-Archimedean criterion}
Let $|\cdot|$ be an absolute value on a field $K$. The following are
equivalent.
\begin{enumerate}
\item $|\cdot|$ is non-Archimedean.
\item $|n\cdot1_K|\leq1$ for every $n\in\mathbf Z$.
\item The set $\{|n\cdot1_K|\mid n\in\mathbf Z\}$ is bounded.
\end{enumerate}
\end{lemma}
\begin{proof}
(1)$\Rightarrow$(2). By induction,
$|n\cdot1|=|1+\cdots+1|\leq\max\{|1|,\ldots,|1|\}=1$ for $n\geq0$, and
$|-a|=|a|$ because $|-1|^2=1$ forces $|-1|=1$.

(2)$\Rightarrow$(3) is trivial.

(3)$\Rightarrow$(1). Let $M$ bound all $|n\cdot1|$ and let $a,b\in K$. For
every $N\geq1$, the binomial theorem and the triangle inequality give
\[
 |a+b|^N=\left|\sum_{i=0}^N\binom{N}{i}a^ib^{N-i}\right|
 \leq\sum_{i=0}^{N}M|a|^i|b|^{N-i}
 \leq (N+1)M\max\{|a|,|b|\}^N.
\]
Taking $N$-th roots and letting $N\to\infty$ yields
$|a+b|\leq\max\{|a|,|b|\}$, since $((N+1)M)^{1/N}\to1$.
\end{proof}

\begin{example}\label{ex:basic-absolute-values}
\index{absolute value!examples}
(1) The trivial absolute value exists on every field; its valuation is
$v(a)=0$ for $a\ne0$. It is complete, since every Cauchy sequence is
eventually constant.

(2) For a prime $p$, write a nonzero rational number as
$a=p^{m}u$ with $u$ a quotient of integers prime to $p$; then
$|a|_p:=p^{-m}$ is a non-Archimedean absolute value on $\mathbf Q$ by
\cref{lem:ultrametric-integer-bound}, with valuation
$v_p(a)=m\log p$. The frequent normalization $v_p(p)=1$ replaces the base
$\rme$ by the base $p$; by the remark above, this changes nothing
structural.

(3) On the Laurent series field $k((t))$, the assignment
$v(f)=\operatorname{ord}_t(f)$ is a complete discretely valued absolute
value $|f|=\rme^{-\operatorname{ord}_t f}$, trivial on $k$. Restricting
along $k(t)\subset k((t))$ gives the $t$-adic valuation of the rational
function field. This pair is the algebraic model for all
equicharacteristic discretely valued fields appearing in our
degeneration pictures.
\end{example}

\begin{lemma}\label{lem:strong-triangle-equality}
If $v(f)\ne v(g)$, then
\[
                           v(f+g)=\min\{v(f),v(g)\}.
\]
\end{lemma}
\begin{proof}
Suppose $v(f)<v(g)$. The ultrametric inequality gives
$v(f+g)\geq v(f)$. Since $f=(f+g)-g$, it also gives
\[
 v(f)\geq\min\{v(f+g),v(g)\}.
\]
As $v(g)>v(f)$, the latter minimum can be at most $v(f)$ only if
$v(f+g)\leq v(f)$. Hence equality holds.
\end{proof}

The proof uses only the ordered structure of the value set. The lemma and
its proof therefore apply verbatim to valuations with values in an
arbitrary totally ordered abelian group, as introduced later in this
chapter, and we will use it in that generality without further comment. A
useful reformulation: in a sum of finitely many terms with pairwise
comparisons available, if exactly one term has strictly smallest value,
then the sum is nonzero and has that value. Most nonvanishing arguments in
this chapter are instances of this reformulation.

\begin{definition}\label{def:semivaluation-support}
The \emph{support} of a semivaluation $v$ is
\[
                         \operatorname{supp}(v)
          =\{f\in A\mid v(f)=+\infty\}.
\]
\end{definition}

\begin{proposition}
\label{prop:semivaluation-support-prime}
The support of a semivaluation is a prime ideal. Moreover, $v$ induces a
valuation on the fraction field of $A/\operatorname{supp}(v)$.
\end{proposition}
\begin{proof}
The support contains $0$ and not $1$. It is closed under addition by the
ultrametric inequality and under multiplication by arbitrary elements by
multiplicativity. If $fg$ lies in the support, then
$v(f)+v(g)=+\infty$, so at least one of $v(f),v(g)$ is $+\infty$.
Thus the support is prime.

First the value descends to $A/\operatorname{supp}(v)$.  Indeed, if
$f-f'$ belongs to the support and one of $v(f),v(f')$ is finite, the strong
triangle equality applied to $f=f'+(f-f')$ gives $v(f)=v(f')$; if one is
infinite, then both representatives lie in the support.  All nonzero classes
therefore have finite value.  Define $v(f/g)=v(f)-v(g)$. If
$f/g=f'/g'$, then $fg'=f'g$ and multiplicativity shows that the definition
is independent of the presentation. The valuation axioms extend immediately.
\end{proof}

\begin{example}\label{ex:trivial-point-evaluation-semivaluation}
Evaluation at a point $a$ of an affine $k$-variety gives the semivaluation
$v_a(f)=0$ if $f(a)\ne0$ and $v_a(f)=+\infty$ if $f(a)=0$, when $k$ is
trivially valued. Its support is the prime ideal of the Zariski closure of
$a$. Thus proper-support points cannot be discarded from a trivially valued
analytification.
\end{example}

\section{Monomial and quasimonomial valuations}

\begin{definition}\label{def:monomial-quasimonomial-valuation}
Let $Y$ be a regular variety over a field, let $\xi$ be the generic point
of an integral closed subvariety of codimension $r$, and let
$z_1,\ldots,z_r$ be a regular system of parameters of
$\mathcal O_{Y,\xi}$. Choose
$\alpha=(\alpha_1,\ldots,\alpha_r)\in\mathbf R_{>0}^r$.  After choosing a
coefficient field in the equicharacteristic completion, write
\[
 f=\sum_{\beta\in\mathbf N^r}c_\beta z^\beta,
\]
with $c_\beta\in\kappa(\xi)$.  Equivalently, read the coefficients from
the monomial basis of
$\operatorname{gr}_{(z_1,\ldots,z_r)}\mathcal O_{Y,\xi}
\simeq\kappa(\xi)[\bar z_1,\ldots,\bar z_r]$.  The
associated \emph{monomial valuation} is
\[
                     v_\alpha(f)=
       \min_{c_\beta\ne0}\langle\alpha,\beta\rangle.
\]
A valuation on $K(X)$ is \emph{quasimonomial} if it becomes monomial on some
regular birational model in this way.
\end{definition}

The resulting least weight is independent of the chosen coefficient field:
changing a lift of a residue coefficient changes it only by terms divisible
by at least one $z_i$, hence by terms of strictly larger positive weight.
The minimum exists when all $\alpha_i>0$: below each fixed bound
only finitely many exponent vectors occur. Cancellation can only raise the
value of a sum. For multiplicativity, the weighted associated graded ring is
a polynomial ring over the residue field in the initial forms of the $z_i$,
hence is a domain. Thus lowest-weight parts have nonzero product.

Two remarks orient the reader. First, the definition depends on the model
$Y$, the point $\xi$, and the weights, but the resulting valuation on the
function field is intrinsic: its value group, rank, rational rank, and
residue field are invariants of the valuation alone and can be computed on
any model realizing it. These invariants are computed completely, by way of
graded reduction, in \cref{lem:monomial-invariants} below. Second,
quasimonomial valuations are exactly the points that live on skeleta: on an
SNC model, the monomial valuations in the parameters cutting out a stratum
sweep out a simplex, and these simplices glue into the dual complex. This
is carried out in \cref{chap:skeleta}, and the same construction over a
trivially valued field appears in \cref{chap:trivial-analytification}. When
the weights are rationally independent, the valuation remembers the whole
simplex direction; when they are commensurable, it degenerates towards a
divisorial valuation on a blowup. The gallery in
\cref{sec:valuation-gallery} makes both behaviours explicit.

\begin{example}[Monomial ideals]\label{ex:monomial-ideal-valuation}
For $A=k[x_1,\ldots,x_n]$ and
$\alpha\in\mathbf R_{\geq0}^n$, put
\[
 v_\alpha\left(\sum_\beta c_\beta x^\beta\right)
       =\min_{c_\beta\ne0}\langle\alpha,\beta\rangle.
\]
If $\mathfrak a=(x^{\beta^{(1)}},\ldots,x^{\beta^{(N)}})$, then
\[
                  v_\alpha(\mathfrak a)
       =\min_j\langle\alpha,\beta^{(j)}\rangle.
\]
Thus the valuation of a monomial ideal is the support function, with a
minimum convention, of its Newton polyhedron.
\end{example}

The monomial-ideal computation is the germ of the tropical point of view:
the family of all monomial valuations, evaluated on the semiring of
monomial ideals, is a compact tropical spectrum. This is proved in
\cref{ex:monomial-tropical-spectrum} and generalized in
\cref{app:tropical-algebra} and \cref{chap:tropical-spectra-complex}.

\begin{example}[Valuations on a smooth curve]\label{ex:smooth-curve-valuations}
Let $C$ be a smooth projective curve over a trivially valued algebraically
closed field. If $p\in C$ and $c>0$, then
$c\operatorname{ord}_p$ is a real valuation centered at $p$. Since every
proper birational morphism between smooth projective curves is an
isomorphism, there are no new divisorial valuations on higher smooth models.
Proper-support semivaluations still add the scheme points themselves.
\end{example}

\section{Completed residue fields}

Let $A$ be a Banach $K$-algebra and let $x$ be a bounded multiplicative
seminorm on $A$. Set $\mathfrak p_x=\operatorname{supp}(v_x)$. The induced
absolute value on the fraction field of $A/\mathfrak p_x$ has a completion.

\begin{definition}\label{def:completed-residue-field}
The resulting complete valued field is the \emph{completed residue field} at
$x$ and is denoted $\mathcal H(x)$. The composite
\[
                    A\longrightarrow A/\mathfrak p_x
       \longrightarrow\operatorname{Frac}(A/\mathfrak p_x)
       \longrightarrow\mathcal H(x)
\]
is the canonical evaluation map at $x$.
\end{definition}

\begin{proposition}
\label{prop:completed-residue-universal}
For $f\in A$, the norm of the image of $f$ in $\mathcal H(x)$ equals
$|f(x)|$. The pair consisting of $\mathcal H(x)$ and the evaluation map is
initial among isometric maps from $A/\mathfrak p_x$ to complete valued fields.
\end{proposition}
\begin{proof}
The first assertion holds by construction before completion and completion
is isometric. Any isometric map to a field extends uniquely to the fraction
field and then, by completeness, uniquely to its completion. This gives the
universal property.
\end{proof}

\begin{lemma}\label{lem:completion-invariants}
\index{completed residue field}
Let $F$ be a field with a rank-one valuation $v$ and let $\widehat F$ be
its completion. The inclusion $F\subset\widehat F$ induces equalities of
value groups and an isomorphism of residue fields:
\[
 \Gamma_{\widehat F}=\Gamma_F,
 \qquad
 \kappa(\widehat F)\simeq\kappa(F).
\]
\end{lemma}
\begin{proof}
Let $\widehat a\in\widehat F^\times$ and choose $a\in F$ with
$v(\widehat a-a)>v(\widehat a)$, possible by density. The strong triangle
equality (\cref{lem:strong-triangle-equality}) applied to
$\widehat a=a+(\widehat a-a)$ gives $v(a)=v(\widehat a)$. Hence every value
of $\widehat F$ is a value of $F$, and the value groups agree.

For residue fields, take $\widehat a$ with $v(\widehat a)=0$ and choose
$a\in F$ with $v(\widehat a-a)>0$. Then $v(a)=0$ and $\widehat a$, $a$ have
the same image in $\kappa(\widehat F)$. Thus
$\kappa(F)\to\kappa(\widehat F)$ is surjective; it is injective because
$\mathfrak m_{\widehat F}\cap R_F=\mathfrak m_F$.
\end{proof}

Consequently the two classical invariants of a point
$x\in X^{\mathrm{an}}$---the value group $|\mathcal H(x)^\times|$ and the
residue field $\widetilde{\mathcal H(x)}$---can be computed before
completion, on the fraction field of $A/\mathfrak p_x$ itself. The graded
computations of the next section do exactly this. For points of analytic
curves, these two invariants separate types II and III from each other and
from the pair \{I,\,IV\}: a point of type II enlarges the residue field,
a point of type III enlarges the value group, and a point of type I or IV
enlarges neither.  They do \emph{not} separate I from IV, since
$\mathcal H(x)$ is an immediate extension of $K$ in both cases
(\Cref{prop:berk-type-invariants}); that separation is made by the
condition $\mathcal H(x)=K$, which is not a statement about the two
invariants.

Completed residue fields allow analytic points to be used as genuine base
fields. For a morphism $Y\to X$ and $x\in X^{\mathrm{an}}$, the analytic
fibre over $x$ is obtained by base change to $\mathcal H(x)$, not merely to
the ordinary residue field of the support. The extension theory needed to
control such base changes---in particular the fact that a finite extension
of a complete rank-one field carries a unique compatible absolute
value---is developed in \cref{sec:extension-of-valuations}.

\section{Graded reduction}

An ordinary residue field remembers which elements have value zero and
forgets all other leading terms. Graded reduction remembers a leading term
in every value.

\begin{definition}\label{def:valuation-associated-graded}
Let $v:A\to\Gamma\cup\{+\infty\}$ be a semivaluation. For
$\gamma\in\Gamma$, set
\[
 F^{\geq\gamma}A=\{f\in A\mid v(f)\geq\gamma\},\qquad
 F^{>\gamma}A=\{f\in A\mid v(f)>\gamma\}.
\]
The \emph{associated graded ring} is
\[
             \operatorname{gr}_v A=
       \bigoplus_{\gamma\in\Gamma}
       F^{\geq\gamma}A/F^{>\gamma}A.
\]
For $f$ of finite value, its image in the component of degree $v(f)$ is its
\emph{initial form} $\operatorname{in}_v(f)$.
\end{definition}

\begin{proposition}
\label{prop:graded-domain}
If $\operatorname{supp}(v)=0$, then $\operatorname{gr}_v A$ is a domain and
\[
       \operatorname{in}_v(fg)
       =\operatorname{in}_v(f)\operatorname{in}_v(g)\ne0
\]
for nonzero $f,g\in A$.
\end{proposition}
\begin{proof}
Multiplicativity gives $v(fg)=v(f)+v(g)$, so multiplication of the two
classes is the class of $fg$ in degree $v(f)+v(g)$. This class cannot vanish,
for vanishing would mean $v(fg)>v(f)+v(g)$. Every nonzero homogeneous element
is the initial form of an element of $A$; products of nonzero homogeneous
elements are therefore nonzero. Ordering the finitely many degrees appearing
in two arbitrary graded elements and taking their least terms shows that
their product is nonzero.
\end{proof}

The next lemma explains why graded reduction computes the two classical
invariants of a valuation, its value group and its residue field, in one
stroke. Call a graded ring in which every nonzero homogeneous element is
invertible a \emph{graded field}, and for a graded domain $G$ write
$G^{\mathrm{gr}}$ for its homogeneous fraction ring, the localization at
all nonzero homogeneous elements.

\begin{lemma}\label{lem:graded-residue-field}
\index{graded reduction!residue field and value group}
Let $F$ be a field with a valuation $v$, with values in a totally ordered
abelian group, and let $A\subset F$ be a subring with
$\operatorname{Frac}(A)=F$. Then:
\begin{enumerate}
\item $\operatorname{gr}_vF$ is a graded field, and the natural map
      $(\operatorname{gr}_vA)^{\mathrm{gr}}\to\operatorname{gr}_vF$ is an
      isomorphism of graded rings;
\item the degree-zero component of $\operatorname{gr}_vF$ is the residue
      field $\kappa(v)$;
\item the set of degrees in which $\operatorname{gr}_vF$ is nonzero is the
      value group $\Gamma_v$, which is generated as a group by
      $v(A\setminus\{0\})$.
\end{enumerate}
\end{lemma}
\begin{proof}
(1) For $f\in F^\times$, multiplicativity of initial forms gives
$\operatorname{in}_v(f)\operatorname{in}_v(f^{-1})
=\operatorname{in}_v(1)=1$, so every initial form is invertible. A nonzero
homogeneous element of degree $\gamma$ in $\operatorname{gr}_vF$ is the
class of some $f$ with $v(f)=\gamma$, hence is an initial form. Thus
$\operatorname{gr}_vF$ is a graded field.

For each $\gamma$, the natural map
$\{f\in A\mid v(f)\geq\gamma\}/\{f\in A\mid v(f)>\gamma\}
\to F^{\geq\gamma}/F^{>\gamma}$ is injective, so
$\operatorname{gr}_vA\to\operatorname{gr}_vF$ is an injective map of graded
domains. Since $\operatorname{gr}_vF$ is a graded field, the map extends to
the homogeneous fraction ring. Surjectivity: every element of $F^\times$ is
$f/g$ with $f,g\in A\setminus\{0\}$, and multiplicativity gives
$\operatorname{in}_v(f/g)
=\operatorname{in}_v(f)\operatorname{in}_v(g)^{-1}$, a homogeneous
fraction from $A$.

(2) By definition, the degree-zero component of $\operatorname{gr}_vF$ is
$F^{\geq0}/F^{>0}=R_v/\mathfrak m_v=\kappa(v)$.

(3) A degree $\gamma$ occurs iff some $f\in F^\times$ has $v(f)=\gamma$,
which is the definition of $\Gamma_v$. Writing $f=g/h$ with $g,h\in A$
exhibits $\gamma=v(g)-v(h)$ as a difference of values on $A$.
\end{proof}

In practice the lemma is used as follows: one computes
$\operatorname{gr}_vA$ for a convenient finitely generated $A$---often a
polynomial ring, where the graded ring is again a polynomial ring with
shifted degrees---and reads off $\Gamma_v$ as the group generated by the
degrees and $\kappa(v)$ as the field of degree-zero homogeneous fractions.
The computations in \cref{ex:graded-gauss} and
\cref{lem:monomial-invariants} are all of this shape.

For a valued field $K$, the degree-zero fraction of this construction is the
usual residue field, while all degrees together form Temkin's graded residue
field. In multiplicative notation one indexes by radii $r>0$ and uses
$K_{\leq r}/K_{<r}$. The conversion is $r=\rme^{-\gamma}$. Graded reduction is
the natural coefficient object for tropical initial degenerations and for
the tropical spectra used later in the complex-analytic construction; see
\cite{MP24}.

\begin{example}[Initial forms at a Gauss point]\label{ex:gauss-initial-forms}
For the Gauss valuation
$v_r(\sum a_iT^i)=\min_i(v(a_i)+ir)$, the initial form is the sum of precisely
those monomials attaining the minimum. A bend of the function
$r\mapsto v_r(f)$ occurs when at least two monomials have the same initial
value. This elementary observation is the bridge between valuations and
tropical hypersurfaces.
\end{example}

\begin{example}[Graded reduction at Gauss points]\label{ex:graded-gauss}
\index{graded reduction!of a Gauss point}\index{Gauss valuation}
Let $K$ be a complete discretely valued field, normalized by
$v_K(\pi)=1$, with residue field $\widetilde K$. We compute the graded
reduction of the Gauss valuations
$v_r(\sum_ia_iT^i)=\min_i(v_K(a_i)+ir)$ on $A=K[T]$, for $r>0$, and read
off their invariants by \cref{lem:graded-residue-field}.

\emph{Step 1: the graded reduction of $K$.} The component of
$\operatorname{gr}_{v_K}K$ in degree $\gamma$ is
$K_{\geq\gamma}/K_{>\gamma}$, which vanishes unless $\gamma\in\mathbf Z$
and is then isomorphic to $\widetilde K$ via $a\mapsto\overline{a/\pi^m}$
for $\gamma=m$. Writing $S=\operatorname{in}_{v_K}(\pi)$, in degree $1$,
multiplicativity of initial forms gives
$S\cdot\operatorname{in}(\pi^{-1})=1$, so
\[
 \operatorname{gr}_{v_K}K=\widetilde K[S,S^{-1}],
\]
a Laurent polynomial ring graded by $\deg S=1$.

\emph{Step 2: the graded reduction of $K[T]$.} Because
$v_r(\sum a_iT^i)=\min_i(v_K(a_i)+ir)$ is computed coefficientwise,
\[
 F^{\geq\gamma}A=\bigoplus_{i\geq0}K_{\geq\gamma-ir}\,T^i,
 \qquad
 F^{>\gamma}A=\bigoplus_{i\geq0}K_{>\gamma-ir}\,T^i,
\]
so the degree-$\gamma$ component of $\operatorname{gr}_{v_r}A$ is
$\bigoplus_i(K_{\geq\gamma-ir}/K_{>\gamma-ir})$. Setting
$\Theta=\operatorname{in}_{v_r}(T)$, in degree $r$, we conclude that
\[
 \operatorname{gr}_{v_r}K[T]
 =(\operatorname{gr}_{v_K}K)[\Theta]
 =\widetilde K[S,S^{-1}][\Theta],
\]
a polynomial ring over $\widetilde K[S^{\pm1}]$ with $\deg\Theta=r$; the
ring structure agrees because initial forms multiply. In particular the
graded ring is a domain, which reproves multiplicativity of $v_r$, that is,
the Gauss lemma.

\emph{Step 3: invariants.} By \cref{lem:graded-residue-field}, the value
group of $v_r$ on $K(T)$ is generated by the degrees:
$\Gamma_{v_r}=\mathbf Z+r\mathbf Z$. The residue field is the field of
degree-zero homogeneous fractions of $\widetilde K[S^{\pm1},\Theta]$, and a
homogeneous fraction of monomials $S^m\Theta^n$ has degree $m+nr$. Two
cases appear.
\begin{enumerate}
\item \emph{$r=a/b$ rational, in lowest terms with $b>0$.} Then
      $\Gamma_{v_r}=\tfrac1b\mathbf Z$, so the value group index is
      $[\Gamma_{v_r}:\Gamma_K]=b$. The lattice of pairs $(m,n)$ with
      $m+nr=0$ is generated by $(-a,b)$, so the degree-zero part is
      generated over $\widetilde K$ by $u:=S^{-a}\Theta^b$, the class of
      $T^b/\pi^a$. Distinct powers of $u$ are linearly independent over
      $\widetilde K$ inside the graded field, so $u$ is transcendental
      and
      \[
       \kappa(v_r)=\widetilde K(u),
       \qquad
       \operatorname{trdeg}_{\widetilde K}\kappa(v_r)=1.
      \]
\item \emph{$r$ irrational.} Then $m+nr=0$ forces $m=n=0$: the value
      group is $\Gamma_{v_r}=\mathbf Z\oplus r\mathbf Z$, free of rank
      two, and
      \[
       \kappa(v_r)=\widetilde K,
       \qquad
       \operatorname{trdeg}_{\widetilde K}\kappa(v_r)=0.
      \]
      Every $f$ has a single monomial as initial form, because the
      degrees $v_K(a_i)+ir$ are pairwise distinct.
\end{enumerate}
In both cases exactly one of the two invariants grows, and
\[
 \dim_{\mathbf Q}
 \bigl((\Gamma_{v_r}/\Gamma_K)\otimes_{\mathbf Z}\mathbf Q\bigr)
 +\operatorname{trdeg}_{\widetilde K}\kappa(v_r)
 =1
 =\operatorname{trdeg}_KK(T),
\]
so Gauss points saturate the relative Abhyankar inequality proved in
\cref{thm:abhyankar-relative} below. After base change to an algebraically
closed field, the two cases correspond to points of type II and type III on
the analytic line in the classification of \cref{sec:berk-line}.
\end{example}

\begin{example}[The tropical spectrum of monomial ideals]
\label{ex:monomial-tropical-spectrum}
\index{tropical spectrum!of monomial ideals}
Graded reduction records one leading term per value; tropical spectra
record the values themselves, on the idempotent semiring of ideals. We
compute one tropical spectrum in full. Let $k$ be a field and let
$\mathscr M$ be the set of monomial ideals of $k[x,y]$, including $(0)$ and
$(1)$, with the operations
\[
 I\oplus J=I+J,\qquad I\odot J=IJ.
\]
Sums and products of monomial ideals are monomial, $I+I=I$, and both
operations are commutative, associative, and distributive; thus
$\mathscr M$ is an idempotent semiring, a subsemiring of the ideal
semiring of \cref{def:appG-ideal-semiring}. Following
\cref{def:appG-tropical-spectrum} and \cref{chap:tropical-spectra-complex},
its tropical spectrum is the set of semiring homomorphisms
\[
 \chi:\mathscr M\longrightarrow([0,+\infty],\min,+),
 \qquad
 \chi((0))=+\infty,\quad\chi((1))=0,
\]
with the topology of pointwise convergence. We claim that evaluation at the
two coordinate ideals is a homeomorphism
\[
 \operatorname{tSpec}(\mathscr M)\xrightarrow{\ \sim\ }[0,+\infty]^2,
 \qquad
 \chi\longmapsto(\chi((x)),\chi((y))).
\]

\emph{The map is injective.} First,
$(x^iy^j)=(x)^i(y)^j$, so multiplicativity gives
$\chi((x^iy^j))=i\chi((x))+j\chi((y))$, where a summand with zero
exponent is omitted. Thus, for example, the value of $(y^j)$ is
$j\chi((y))$ even when $\chi((x))=+\infty$; no product
$0\cdot(+\infty)$ is formed. Every nonzero monomial ideal is a finite sum of principal
monomial ideals, so $\chi(I)$ is the minimum of $i\chi((x))+j\chi((y))$
over a finite generating set of monomials $x^iy^j$ of $I$, with the same
zero-exponent omission. Hence $\chi$ is
determined by its values on $(x)$ and $(y)$.

\emph{The map is surjective.} Given $(a,b)\in[0,+\infty]^2$, define
$\chi_{a,b}(I)=\min\{ia+jb\mid x^iy^j\in I\}$ for $I\ne(0)$, and
$\chi_{a,b}((0))=+\infty$, again omitting every zero-exponent summand.
The minimum over all monomials of $I$ equals
the minimum over any monomial generating set: a monomial lies in a
monomial ideal precisely when it is divisible by one of the generators---%
expanding a representation $x^iy^j=\sum f_s\cdot g_s$ with monomial
generators $g_s$ and comparing coefficients of $x^iy^j$ exhibits a
monomial factor---and divisibility can only increase the linear form
$(i,j)\mapsto ia+jb$ since $a,b\geq0$. With this description,
$\chi_{a,b}(I+J)=\min\{\chi_{a,b}(I),\chi_{a,b}(J)\}$ because the union of
two generating sets generates the sum, and
$\chi_{a,b}(IJ)=\chi_{a,b}(I)+\chi_{a,b}(J)$ because the pairwise products
generate the product and the minimum of $(i+i')a+(j+j')b$ over independent
pairs splits as the sum of the two minima. Thus $\chi_{a,b}$ is a
character, and it takes $((x),(y))$ to $(a,b)$.

\emph{Topologies match.} The bijection is continuous in both directions:
coordinates of the image are evaluations at $(x)$ and $(y)$, while the
value at any $I$ is the minimum of finitely many continuous functions
$ia+jb$ of $(a,b)$, using again that a fixed finite generating set
computes it.

The finite points with $a,b<+\infty$ are precisely the monomial
semivaluations of the earlier example in this chapter; the boundary strata
at $+\infty$ are semivaluations with nonzero support, for instance
$a=+\infty$ corresponds to supports containing $(x)$. The square
$[0,+\infty]^2$ is the nonnegative part of the extended tropical affine
plane of \cref{app:tropical-algebra}. This baby case displays the general
mechanism of \cref{thm:appG-ideal-characters}: characters of an ideal
semiring are semivaluations, and compactness of the spectrum comes from
the product topology on a cube of coordinates, exactly as for the
tropical spectra of complex spaces in
\cref{chap:tropical-spectra-complex}, following \cite{MP24}.
\end{example}

\section{Valuation rings and centers}

We briefly allow values in a totally ordered abelian group $\Gamma$. A
valuation on a field $F$ is a homomorphism
$v:F^\times\to\Gamma$ satisfying the ultrametric inequality, extended by
$v(0)=+\infty$.

\begin{definition}\label{def:valuation-ring-residue-field}
The \emph{valuation ring}, maximal ideal, and residue field of $v$ are
\[
 R_v=\{f\in F\mid v(f)\geq0\},\qquad
 \mathfrak m_v=\{f\in F\mid v(f)>0\},\qquad
 \kappa(v)=R_v/\mathfrak m_v.
\]
Its value group is $\Gamma_v=v(F^\times)$.
\end{definition}

\begin{proposition}\label{prop:valuation-ring-characterization}
For every $f\in F^\times$, either $f\in R_v$ or $f^{-1}\in R_v$. Conversely,
if a subring $R\subset F$ has this property, then $R$ is a valuation ring of
$F$ for a valuation with values in a totally ordered abelian group.
\end{proposition}
\begin{proof}
The first assertion follows because either $v(f)\geq0$ or
$v(f^{-1})=-v(f)\geq0$. For the converse, the group
$\Gamma=F^\times/R^\times$ becomes totally ordered by
\[
       [f]\leq[g]\quad\Longleftrightarrow\quad g/f\in R.
\]
The assumed alternative gives totality, and a direct check gives a
translation-invariant order. The quotient map $F^\times\to\Gamma$, with the
order reversed if necessary, is a valuation whose nonnegative elements are
exactly $R$. See \cite[Chapter~VI, \S1]{ZariskiSamuelII} for the full
ordered-group formulation.
\end{proof}

\begin{definition}\label{def:valuation-center}
Let $X$ be an integral scheme with function field $F$. A valuation $v$ of
$F$ has \emph{center} $\xi\in X$ if
\[
                  \mathcal O_{X,\xi}\subset R_v,
             \qquad \mathfrak m_\xi=\mathfrak m_v\cap\mathcal O_{X,\xi}.
\]
Equivalently, the generic-point morphism $\operatorname{Spec}F\to X$ extends
to a morphism $\operatorname{Spec}R_v\to X$ carrying the closed point to
$\xi$.
\end{definition}

\begin{proposition}\label{prop:centers-proper-separated}
Let $X$ be an integral scheme with function field $F$, of finite type over a
base scheme $S$, and let $v$ be a valuation of $F$ for which the composite
$\operatorname{Spec}F\to X\to S$ extends to a morphism
$\operatorname{Spec}R_v\to S$.  If $X$ is separated over $S$, then $v$ has at
most one center on $X$; if $X$ is proper over $S$, it has exactly one.  In
particular a valuation of $F$ trivial on $k$ has a unique center on a proper
$k$-scheme, and a valuation of $F$ whose ring contains $K^\circ$ has a unique
center on a proper $K^\circ$-model.  The uniqueness half needs no hypothesis
on $v$ beyond separatedness.
\end{proposition}

\begin{warning}{Existence needs the base}
Without the extension hypothesis the center need not exist, even on a proper
scheme.  Take $k=\mathbf C$, $X=\mathbf P^1_{\mathbf C}$ and
$F=\mathbf C(z)$.  Since $\mathbf C$ has infinite transcendence degree over
$\mathbf Q$ it carries nontrivial valuations; extend one to $F$ by
\Cref{prop:chevalley-extension}.  A center $\xi$ would force
$\mathbf C\subset\cO_{X,\xi}\subset R_v$, contradicting nontriviality on
$\mathbf C$.  This is exactly the gap relative to
\cite[Tag 0BX5]{stacks-project}, which is stated for a given morphism
$\operatorname{Spec}R_v\to S$.  In this book the hypothesis always holds,
because the valuations considered are either trivial on the ground field or
normalized so that $R_v\supseteq K^\circ$.
\end{warning}
\begin{proof}
This is the valuative criterion. Two extensions
$\operatorname{Spec}R_v\to X$ agreeing at the generic point agree when the
diagonal of $X$ is closed; this proves uniqueness. Properness, applied to the given morphism
$\operatorname{Spec}R_v\to S$, gives an extension, hence existence. The uniqueness statement is the separatedness
implication in \cite[Tag 01KZ]{stacks-project}, while existence and
uniqueness for a proper morphism are given by
\cite[Tag 0BX5]{stacks-project}.
\end{proof}

Centers are the algebraic shadow of the analytic reduction maps of
\cref{sec:berk-reduction}: for a model of a proper variety, sending an
analytic point to the center of its valuation on the model is exactly the
specialization map to the special fibre. Divisorial points, whose centers
are generic points of components of special fibres, are studied in
\cref{chap:models}.

\section{Extensions of valuations}
\label{sec:extension-of-valuations}
\index{valuation!extension of}

Analytic geometry constantly changes the field: to a finite extension, to
a completed residue field, to a bigger algebraically closed field. This
section records what happens to valuations under such changes. The
existence theorem is due to Chevalley, and the numerical bookkeeping is
the fundamental inequality. Over a complete rank-one field the extension
is unique, which is the algebraic reason that analytic fibres are
well-defined spaces rather than sets of choices.

\begin{proposition}
\label{prop:chevalley-extension}
\index{Chevalley extension theorem}
Let $F\subset L$ be a field extension and let $v$ be a valuation on $F$,
with values in a totally ordered abelian group. Then there is a valuation
$w$ on $L$ with
\[
 R_w\cap F=R_v,
 \qquad
 \mathfrak m_w\cap F=\mathfrak m_v.
\]
Moreover the inclusion $F^\times\subset L^\times$ induces an
order-preserving embedding $\Gamma_v\hookrightarrow\Gamma_w$ and an
embedding of residue fields $\kappa(v)\hookrightarrow\kappa(w)$.
\end{proposition}
\begin{proof}
\emph{Step 1: a maximal dominating local ring.} Consider the set of local
subrings $(\mathcal O,\mathfrak m_{\mathcal O})$ of $L$ dominating
$(R_v,\mathfrak m_v)$, that is, containing $R_v$ and satisfying
$\mathfrak m_{\mathcal O}\cap R_v=\mathfrak m_v$, partially ordered by
domination. The set contains $R_v$. The union of a chain is a subring
$\mathcal O$ with ideal $\mathfrak m$ given by the union of the maximal
ideals; every element outside $\mathfrak m$ lies outside the maximal ideal
of some member of the chain, hence is invertible there and in
$\mathcal O$. Thus the union is again a local ring dominating $R_v$, and
Zorn's lemma provides a maximal element $(\mathcal O,\mathfrak m)$.

\emph{Step 2: $\mathcal O$ is a valuation ring of $L$.} Let $x\in L^\times$.
We claim that $\mathfrak m\mathcal O[x]\neq\mathcal O[x]$ or
$\mathfrak m\mathcal O[x^{-1}]\neq\mathcal O[x^{-1}]$. Otherwise there are
relations
\[
 1=\sum_{i=0}^{m}a_ix^i,
 \qquad
 1=\sum_{j=0}^{n}b_jx^{-j},
 \qquad a_i,b_j\in\mathfrak m,
\]
with $m$ and $n$ chosen minimal, and, exchanging $x$ and $x^{-1}$ if
necessary, $m\geq n$. Since $1-b_0$ is a unit of $\mathcal O$, the second
relation gives
\[
 x^{n}=\sum_{j=1}^{n}\frac{b_j}{1-b_0}\,x^{n-j},
\]
with coefficients in $\mathfrak m$. Substituting
$x^{m}=x^{m-n}\cdot x^{n}$ into the first relation replaces the term
$a_mx^m$ by terms of degree at most $m-1$ with coefficients in
$\mathfrak m$, contradicting minimality of $m$. This proves the claim.

Say $\mathfrak m\mathcal O[x]$ is a proper ideal. Choose a maximal ideal
$M\supset\mathfrak m\mathcal O[x]$ of $\mathcal O[x]$. Then
$M\cap\mathcal O$ is a proper ideal containing $\mathfrak m$, hence equals
$\mathfrak m$, so the local ring $\mathcal O[x]_M$ dominates $\mathcal O$
and, by transitivity, $R_v$. Maximality forces
$\mathcal O[x]_M=\mathcal O$, whence $x\in\mathcal O$. Thus for every
$x\in L^\times$ either $x$ or $x^{-1}$ lies in $\mathcal O$, and
\Cref{prop:valuation-ring-characterization} makes $\mathcal O$ the
valuation ring of a valuation $w$ on $L$.

\emph{Step 3: restriction to $F$.} The ring $\mathcal O\cap F$ is a
valuation ring of $F$ containing $R_v$, with maximal ideal
$\mathfrak m\cap F$ lying over $\mathfrak m_v$; so $\mathcal O\cap F$
dominates $R_v$. A valuation ring admits no proper domination inside its
own field: if $R_v\subsetneq R$ were local rings of $F$ with
$\mathfrak m_R\cap R_v=\mathfrak m_v$ and $x\in R\setminus R_v$, then
$x^{-1}\in\mathfrak m_v\subset\mathfrak m_R$, so
$1=x\cdot x^{-1}\in\mathfrak m_R$, absurd. Hence $\mathcal O\cap F=R_v$
and $\mathfrak m\cap F=\mathfrak m_v$.

Finally, $F^\times\cap\mathcal O^\times=R_v^\times$, so the natural map
$\Gamma_v=F^\times/R_v^\times\to L^\times/\mathcal O^\times=\Gamma_w$ is
an injective homomorphism; it preserves order because nonnegativity means
membership in the valuation ring on both sides. The residue map
$R_v\to\kappa(w)$ has kernel $\mathfrak m_w\cap R_v=\mathfrak m_v$,
embedding $\kappa(v)$ into $\kappa(w)$. This is the existence theorem for
places of \cite[Chapter~VI, \S4]{ZariskiSamuelII}.
\end{proof}

\begin{lemma}\label{lem:algebraic-extension-invariants}
Let $L/F$ be an algebraic extension and let $w$ be a valuation on $L$
extending $v$ on $F$. Then $\Gamma_w/\Gamma_v$ is a torsion group and
$\kappa(w)/\kappa(v)$ is an algebraic extension. Moreover intersection
with $\Gamma_v$ is a bijection between the convex subgroups of $\Gamma_w$
and those of $\Gamma_v$; in particular $w$ and $v$ have the same rank and
the same rational rank.
\end{lemma}
\begin{proof}
Let $x\in L^\times$ and let $\sum_{i\in S}c_ix^i=0$ be a polynomial
relation with all $c_i\in F^\times$ and $S$ a finite set of at least two
exponents; such a relation exists because $x$ is algebraic and nonzero.
If the values $v(c_i)+iw(x)$ were pairwise distinct, the term of least
value would be unique and the sum could not vanish, by the reformulation
of \cref{lem:strong-triangle-equality}. Hence
$v(c_i)+iw(x)=v(c_j)+jw(x)$ for some $i\neq j$, so
$(i-j)w(x)=v(c_j)-v(c_i)\in\Gamma_v$. Thus $\Gamma_w/\Gamma_v$ is
torsion.

Let $\xi\in R_w$ with residue $\bar\xi$, and let
$\sum_{i\in S}c_i\xi^i=0$ with $c_i\in F^\times$ as before. Dividing by a
coefficient of least value we may assume all $c_i\in R_v$ and some
$c_{i_0}=1$. Reducing modulo $\mathfrak m_w$ gives a nonzero polynomial
over $\kappa(v)$ vanishing at $\bar\xi$. Hence $\kappa(w)/\kappa(v)$ is
algebraic, and rational ranks agree because
$\Gamma_v\otimes\mathbf Q\to\Gamma_w\otimes\mathbf Q$ is an isomorphism
for a torsion quotient.

For convex subgroups: if $\Delta\subset\Gamma_w$ is convex and
$\gamma\in\Delta$, choose $n\geq1$ with $n\gamma\in\Gamma_v$; then
$n\gamma\in\Delta\cap\Gamma_v$, and conversely $0\leq|\gamma|\leq|n\gamma|$
recovers $\gamma$ from $\Delta\cap\Gamma_v$ by convexity. Thus $\Delta$ is
determined by its trace, and every convex subgroup of $\Gamma_v$ arises
from its saturation $\{\gamma\in\Gamma_w\mid n\gamma\in\Delta_0
\text{ for some }n\geq1\}$, which is convex. Ranks therefore agree.
\end{proof}

\begin{proposition}
\label{prop:fundamental-inequality}
\index{fundamental inequality}
Let $L/F$ be a finite extension of degree $n$, let $v$ be a valuation on
$F$, and let $w$ be a valuation on $L$ extending $v$. Set
\[
 e=[\Gamma_w:\Gamma_v],
 \qquad
 f=[\kappa(w):\kappa(v)].
\]
Then $ef\leq n$. More precisely, if $a_1,\ldots,a_r\in L^\times$ have
values in pairwise distinct cosets of $\Gamma_v$ and
$b_1,\ldots,b_s\in R_w^\times$ have residues linearly independent over
$\kappa(v)$, then the products $a_ib_j$ are linearly independent over
$F$; hence $rs\leq n$.
\end{proposition}
\begin{proof}
Suppose $\sum_{i,j}c_{ij}a_ib_j=0$ with $c_{ij}\in F$ not all zero.
Fix $i$ with some $c_{ij}\neq0$ and consider
$u_i:=\sum_jc_{ij}b_j$. Let $j_0$ minimize $v(c_{ij})$ among the nonzero
coefficients. Then $c_{ij}/c_{ij_0}\in R_v$ for all $j$, and
\[
 \frac{u_i}{c_{ij_0}}=\sum_j\frac{c_{ij}}{c_{ij_0}}\,b_j
\]
lies in $R_w$ and has residue
$\sum_j\lambda_j\bar b_j$ with $\lambda_j\in\kappa(v)$ and
$\lambda_{j_0}=1$. Linear independence of the residues makes this residue
nonzero, so $w(u_i/c_{ij_0})=0$, that is,
\[
 w(u_i)=v(c_{ij_0})\in\Gamma_v .
\]
Now consider the outer sum $\sum_iu_ia_i=0$. For those $i$ with
$u_i\neq0$, the values
$w(u_ia_i)=w(u_i)+w(a_i)$ lie in the cosets $w(a_i)+\Gamma_v$, which are
pairwise distinct; so the values are pairwise distinct, the least is
attained exactly once, and the sum cannot vanish unless every $u_i$
is zero. But then every row
$\sum_jc_{ij}b_j=0$ with not all $c_{ij}$ zero contradicts, after the
same normalization, the nonvanishing of the residue. Hence all $c_{ij}$
vanish, proving linear independence, and $rs\leq\dim_FL=n$. Taking
suprema over $r\leq e$ and $s\leq f$ gives $ef\leq n$. The finer
statement---the sum $\sum_ie_if_i$ over all extensions $w_i$ of $v$ to
$L$ is at most $n$---is the fundamental inequality of
\cite[Chapter~VI, \S8]{AC}; we do not use it in this book.
\end{proof}

For complete rank-one fields the extension is unique. The key is the
following equivalence of norms, which will also give the completeness of
finite extensions used tacitly whenever a fibre $\mathcal H(x)$ is
enlarged.

\begin{lemma}\label{lem:norm-equivalence-complete}
\index{norm!on a finite-dimensional vector space}
Let $K$ be a field that is complete with respect to a non-Archimedean
absolute value $|\cdot|$, possibly trivial, and let $V$ be a
finite-dimensional $K$-vector space. Then all ultrametric $K$-norms on
$V$ are equivalent, and $V$ is complete with respect to any of them.
Here an ultrametric $K$-norm is a function
$\|\cdot\|:V\to\mathbf R_{\geq0}$ vanishing only at $0$ and satisfying
$\|\lambda x\|=|\lambda|\,\|x\|$ and
$\|x+y\|\leq\max\{\|x\|,\|y\|\}$.
\end{lemma}
\begin{proof}
Fix a basis $e_1,\ldots,e_n$ of $V$ and let
$\|\sum x_ie_i\|_\infty=\max_i|x_i|$ be the associated sup-norm, for
which $V$ is complete because $K$ is. We show by induction on $n$ that
every ultrametric norm $\|\cdot\|$ is equivalent to $\|\cdot\|_\infty$.

For $n=1$ the norm is $\|x_1e_1\|=|x_1|\,\|e_1\|$, proportional to the
sup-norm. Let $n\geq2$ and put $W=\operatorname{span}(e_1,\ldots,e_{n-1})$.
One inequality is general:
$\|\sum x_ie_i\|\leq\max_i|x_i|\|e_i\|\leq C_1\|x\|_\infty$ with
$C_1=\max_i\|e_i\|$.

By induction, $\|\cdot\|$ restricted to $W$ is equivalent to the sup-norm
of $W$, so $(W,\|\cdot\|)$ is complete. We claim
$d:=\inf_{w\in W}\|e_n-w\|>0$. If $d=0$, choose $w_m\in W$ with
$\|e_n-w_m\|\to0$; then
$\|w_m-w_{m'}\|\leq\max\{\|e_n-w_m\|,\|e_n-w_{m'}\|\}\to0$, so $(w_m)$ is
Cauchy in $W$, with limit $w\in W$ satisfying $\|e_n-w\|=0$, that is,
$e_n\in W$, a contradiction.

Now take $x=w+\lambda e_n$ with $w\in W$, $\lambda\in K$. If
$\lambda\neq0$, then
$\|x\|=|\lambda|\,\|e_n+\lambda^{-1}w\|\geq|\lambda|d$, so
$|\lambda|\leq d^{-1}\|x\|$; this holds trivially when $\lambda=0$.
Consequently
\[
 \|w\|\leq\max\{\|x\|,|\lambda|\,\|e_n\|\}
 \leq\bigl(1+d^{-1}\|e_n\|\bigr)\|x\|,
\]
and the induction hypothesis bounds $\|w\|_{\infty}\leq C_2\|w\|$. Hence
\[
 \|x\|_\infty=\max\{\|w\|_\infty,|\lambda|\}
 \leq C_3\|x\|
\]
for a constant $C_3$ depending only on the norm and the basis. The two
displayed inequalities give equivalence, and completeness follows from
completeness in the sup-norm.
\end{proof}

\begin{theorem}\label{thm:complete-field-extension}
\index{valuation!extension to a finite extension}
Let $K$ be complete with respect to a non-Archimedean absolute value
$|\cdot|$ and let $L/K$ be a finite extension of degree $n$. Then
$|\cdot|$ extends to an absolute value on $L$; the extension is unique;
$L$ is complete; and the extension is given by
\[
 |a|_L=\bigl|N_{L/K}(a)\bigr|^{1/n},
 \qquad a\in L.
\]
\end{theorem}
\begin{proof}
\emph{Existence.} Let $v=-\log|\cdot|$ be the associated real valuation.
\Cref{prop:chevalley-extension} provides a valuation $w_0$ on $L$, with
values in some ordered group, whose ring contracts to $R_v$. By
\cref{lem:algebraic-extension-invariants}, $\Gamma_{w_0}/\Gamma_v$ is
torsion, so the assignment
$\gamma\mapsto\frac1m(m\gamma)$, for any $m\geq1$ with
$m\gamma\in\Gamma_v\subset\mathbf R$, is a well-defined order-preserving
embedding $\Gamma_{w_0}\hookrightarrow\mathbf R$ restricting to the
identity on $\Gamma_v$; well-definedness follows since two choices of
$m$ agree after multiplication. Composing $w_0$ with this embedding gives
a real valuation $w$ on $L$ extending $v$, hence an absolute value
$|\cdot|_1=\rme^{-w}$ extending $|\cdot|$.

\emph{Uniqueness.} Let $|\cdot|_1$ and $|\cdot|_2$ be two extensions.
Each is an ultrametric $K$-norm on the $n$-dimensional space $L$, so by
\cref{lem:norm-equivalence-complete} there is $C\geq1$ with
$C^{-1}|a|_2\leq|a|_1\leq C|a|_2$ for all $a$. Applying this to $a^m$ and
using multiplicativity,
\[
 |a|_1=\bigl(|a^m|_1\bigr)^{1/m}
 \leq C^{1/m}\bigl(|a^m|_2\bigr)^{1/m}=C^{1/m}|a|_2
 \xrightarrow[m\to\infty]{}|a|_2,
\]
and symmetrically $|a|_2\leq|a|_1$. Hence the extension is unique, and
$L$ is complete by the same lemma.

\emph{The norm formula.} Let $M$ be a normal closure of $L/K$; it is a
finite extension of $K$, so it carries a unique extension $|\cdot|_M$,
and $|\cdot|_M$ restricts on $L$ to the unique extension $|\cdot|_L$.
For $\sigma\in\operatorname{Aut}(M/K)$, the function
$|\sigma(\cdot)|_M$ is again an extension of $|\cdot|$ to $M$, so
$|\sigma a|_M=|a|_M$ for all $a\in M$. Every $K$-embedding
$\tau:L\to M$ extends to a $K$-automorphism of the normal extension $M$,
whence $|\tau a|_M=|a|_M$ for $a\in L$. The field norm is
\[
 N_{L/K}(a)=\Bigl(\prod_{\tau}\tau(a)\Bigr)^{[L:K]_{\mathrm i}},
\]
the product over the distinct $K$-embeddings $\tau:L\to M$, of which
there are the separable degree $[L:K]_{\mathrm s}$, with
$[L:K]_{\mathrm s}[L:K]_{\mathrm i}=n$. Taking absolute values,
\[
 \bigl|N_{L/K}(a)\bigr|
 =\prod_\tau|\tau a|_M^{[L:K]_{\mathrm i}}
 =\bigl(|a|_L\bigr)^{n},
\]
which is the displayed formula. For the classical treatment of these
statements see \cite[Chapter~VI, \S8]{AC}.
\end{proof}

\begin{remark}\label{rem:complete-field-valuation-extensions}
Two consequences are used constantly. First, the absolute value of a
complete field $K$ extends uniquely to any algebraic extension, in
particular to an algebraic closure, by applying the theorem to each
finite subextension; the completion of the algebraic closure is again a
complete valued field. Second, for a point $x$ of an analytic space, any
finite extension of $\mathcal H(x)$ carries a canonical complete absolute
value; this is what makes fibres of finite morphisms over $x$
well-defined analytic spaces over $\mathcal H(x)$, as used from
\cref{chap:berkovich} onwards. When $K$ is trivially valued and $L/K$ is
algebraic, the theorem also shows that the trivial valuation is the only
extension: this is the reason closed points of a variety over a trivially
valued field carry no moduli of semivaluations beyond scaling of their
supports.
\end{remark}

\section{Rank and Abhyankar invariants}

\begin{definition}\label{def:valuation-rank-invariants}
The \emph{rank} of $v$ is the Krull dimension of $R_v$, equivalently the
length of the longest chain of proper convex subgroups of $\Gamma_v$. Its
\emph{rational rank} is
\[
                    \operatorname{rat.rk}(v)
       =\dim_{\mathbf Q}(\Gamma_v\otimes_{\mathbf Z}\mathbf Q).
\]
For an extension $F/k$ on which $v$ is trivial on $k$, its residue
transcendence degree is
\[
                    \operatorname{trdeg}(v)
       =\operatorname{trdeg}_k\kappa(v).
\]
\end{definition}

Every nontrivial real valuation has rank one, but it may have rational rank
larger than one: a subgroup of $\mathbf R$ can have arbitrarily large
$\mathbf Q$-dimension. Rank and rational rank must therefore not be
interchanged.

The two ranks are nevertheless comparable, and the comparison identifies
rank one with real-valuedness. Recall that a totally ordered abelian group
$\Gamma$ is \emph{archimedean} if for all $\gamma,\delta>0$ there is
$n\geq1$ with $\delta\leq n\gamma$.

\begin{lemma}\label{lem:rank-leq-rational-rank}
For every totally ordered abelian group $\Gamma$,
\[
 \operatorname{rank}(\Gamma)\leq
 \dim_{\mathbf Q}(\Gamma\otimes_{\mathbf Z}\mathbf Q),
\]
where the rank is the length of the longest chain of proper convex
subgroups. In particular
$\operatorname{rank}(v)\leq\operatorname{rat.rk}(v)$ for every valuation.
\end{lemma}
\begin{proof}
Let $0=\Delta_0\subsetneq\Delta_1\subsetneq\cdots\subsetneq\Delta_r=\Gamma$
be a chain of convex subgroups and choose
$\gamma_i\in\Delta_i\setminus\Delta_{i-1}$ for $1\leq i\leq r$. We claim
the $\gamma_i$ are $\mathbf Q$-linearly independent. A convex subgroup
$\Delta$ is saturated: if $n\gamma\in\Delta$ with $n\neq0$, then
$0\leq|\gamma|\leq|n\gamma|$ forces $\gamma\in\Delta$. Suppose
$\sum_in_i\gamma_i=0$ with integers $n_i$ not all zero, and let $i_\ast$
be the largest index with $n_{i_\ast}\neq0$. Then
\[
 n_{i_\ast}\gamma_{i_\ast}=-\sum_{i<i_\ast}n_i\gamma_i
 \in\Delta_{i_\ast-1},
\]
so saturation gives $\gamma_{i_\ast}\in\Delta_{i_\ast-1}$, a
contradiction. Hence $r\leq\dim_{\mathbf Q}(\Gamma\otimes\mathbf Q)$.
\end{proof}

\begin{proposition}\label{prop:rank-one-real-embedding}
\index{rank!one and real values}\index{Holder theorem@H\"older theorem}
A nonzero totally ordered abelian group $\Gamma$ has rank one if and only
if it is archimedean, if and only if it admits an order-preserving
embedding into $(\mathbf R,+)$. Consequently, a nontrivial valuation
admits an equivalent real-valued normalization exactly when its rank is
one.
\end{proposition}
\begin{proof}
\emph{Rank one implies archimedean.} If $\gamma,\delta>0$ and
$\delta>n\gamma$ for all $n$, then
$\Delta=\{\beta\mid |\beta|\leq n\gamma\text{ for some }n\}$ is a convex
subgroup containing $\gamma$ but not $\delta$: it is nonzero and proper,
so the rank is at least two.

\emph{Archimedean implies embeddable.} Fix $\alpha>0$ and define, for
$\gamma\in\Gamma$,
\[
 \varphi(\gamma)
 =\sup\Bigl\{\tfrac pq\in\mathbf Q\ \Bigm|\ q\geq1,\ p\alpha\leq q\gamma
 \Bigr\}.
\]
The defining set is nonempty and bounded above: by the archimedean
property there is $N$ with $|\gamma|\leq N\alpha$, so $-N$ belongs to the
set and every member is at most $N$. Thus
$\varphi:\Gamma\to\mathbf R$ is well-defined.

$\varphi$ is superadditive: from $p\alpha\leq q\gamma$ and
$p'\alpha\leq q\delta$ we get $(p+p')\alpha\leq q(\gamma+\delta)$, and
rationals with common denominators are cofinal in both defining sets.
Conversely, suppose
$\varphi(\gamma+\delta)>\varphi(\gamma)+\varphi(\delta)$. Choose rationals
$\tfrac pq>\varphi(\gamma)$ and $\tfrac{p'}q>\varphi(\delta)$, with a
common denominator $q\geq1$, such that
$\tfrac{p+p'}q<\varphi(\gamma+\delta)$. Since
$\tfrac pq$ exceeds the supremum, $p\alpha\leq q\gamma$ fails, so
$p\alpha>q\gamma$ by totality; likewise $p'\alpha>q\delta$. Adding gives
$(p+p')\alpha>q(\gamma+\delta)$. Now let $\tfrac uw$, with $w\geq1$, be
any member of the set defining $\varphi(\gamma+\delta)$, so that
$u\alpha\leq w(\gamma+\delta)$. Multiplying this inequality by $q$ and the
previous one by $w$ and comparing,
\[
 uq\,\alpha\leq wq(\gamma+\delta)<w(p+p')\,\alpha ,
\]
whence $(uq-w(p+p'))\alpha<0$ and therefore $uq<w(p+p')$, that is,
$\tfrac uw<\tfrac{p+p'}q$. Taking the supremum over $\tfrac uw$ yields
$\varphi(\gamma+\delta)\leq\tfrac{p+p'}q$, contradicting the choice of
$p,p'$. So $\varphi$ is additive.

$\varphi$ is order-preserving with trivial kernel: if $\gamma>0$, the
archimedean property yields $n$ with $\alpha\leq n\gamma$, so
$\varphi(\gamma)\geq\tfrac1n>0$; additivity then gives strict
monotonicity and injectivity.

\emph{Embeddable implies rank one.} A subgroup of $\mathbf R$ is
archimedean, and an archimedean group has no nonzero proper convex
subgroup: if $\Delta\ni\gamma>0$ is convex and $\delta>0$ is arbitrary,
then $\delta\leq n\gamma$ places $\delta$ in $\Delta$.

The final statement follows by composing a valuation of rank one with the
embedding $\varphi$; two such normalizations differ by the choice of
$\alpha$, that is, by a positive scaling. See
\cite[Chapter~VI, \S10]{ZariskiSamuelII}.
\end{proof}

The proposition delimits the reach of Berkovich geometry: the points of
the analytic spaces in this book are real semivaluations, hence valuations
of rank at most one on their support. Higher-rank valuations are not
individually visible as points, but they organize limits and boundaries of
the rank-one locus; \cref{ex:gallery-lex} makes this concrete, and the
Riemann--Zariski perspective is developed in \cite{Tem11}.

\begin{theorem}\label{thm:abhyankar-inequality}
Let $F/k$ be a finitely generated field extension of transcendence degree
$n$, and let $v$ be a valuation of $F$ trivial on $k$. Then
\[
       \operatorname{rat.rk}(v)+\operatorname{trdeg}(v)\leq n.
\]
\end{theorem}
\begin{proof}
Choose elements $f_1,\ldots,f_r\in F^\times$ whose values
$\gamma_i=v(f_i)$ are linearly independent over $\mathbf Q$ modulo the value
group of $k$.  Since $v$ is trivial on $k$, the latter group is zero.  Choose
units $g_1,\ldots,g_s\in R_v^\times$ whose residues
$\overline g_1,\ldots,\overline g_s$ are algebraically independent over
$k\subset\kappa(v)$.  We claim that
\[
                  f_1,\ldots,f_r,g_1,\ldots,g_s
\]
are algebraically independent over $k$.

Indeed, write a purported nonzero polynomial relation as
\[
                  \sum_{a\in\mathbf N^r} f^a P_a(g)=0,
\]
with only finitely many nonzero polynomials
$P_a\in k[T_1,\ldots,T_s]$.  Algebraic independence of the residues implies
$P_a(\overline g)\ne0$ whenever $P_a\ne0$.  Hence
$v(P_a(g))=0$.  The nonzero summands have values
$\langle a,\gamma\rangle$, and these values are pairwise distinct because
the $\gamma_i$ are rationally independent.  The summand of least value is
therefore unique.  The strong triangle equality
\Cref{lem:strong-triangle-equality} says that a sum with a unique term of
least value cannot vanish, a contradiction.

Thus every finite rationally independent family of values, together with
every finite algebraically independent family of residues, has cardinality
at most $\operatorname{trdeg}_kF=n$.  Taking families of maximal possible
sizes gives
\[
 \operatorname{rat.rk}(v)+\operatorname{trdeg}(v)\le n.
\]
This is the standard independence proof of
\cite[Chapter~VI, \S10]{ZariskiSamuelII}.
\end{proof}

The same independence argument proves the version relative to a valued
ground field, which we need for Gauss points and, later, for skeleta over
nontrivially valued fields. The equality case carries finiteness
information that we quote with a precise source.

\begin{theorem}
\label{thm:abhyankar-relative}
\index{Abhyankar inequality!relative form}
Let $L/K$ be a finitely generated field extension of transcendence degree
$n$, let $w$ be a valuation on $L$, and let $v=w|_K$. Then
\[
 \dim_{\mathbf Q}\bigl((\Gamma_w/\Gamma_v)\otimes_{\mathbf Z}\mathbf Q\bigr)
 +\operatorname{trdeg}_{\kappa(v)}\kappa(w)\leq n .
\]
If equality holds, then $\Gamma_w/\Gamma_v$ is a finitely generated
abelian group and $\kappa(w)$ is a finitely generated field extension of
$\kappa(v)$.
\end{theorem}
\begin{proof}
Choose $f_1,\ldots,f_r\in L^\times$ whose values are
$\mathbf Q$-linearly independent modulo $\Gamma_v$, and units
$g_1,\ldots,g_s\in R_w^\times$ whose residues are algebraically
independent over $\kappa(v)$. We claim that
$f_1,\ldots,f_r,g_1,\ldots,g_s$ are algebraically independent over $K$;
this gives $r+s\leq n$ and hence the inequality.

Suppose there is a nontrivial polynomial relation. Grouping monomials by
their exponents in the $f_i$, write it as
\[
 \sum_{a\in\mathbf N^r}f^aP_a(g)=0,
 \qquad
 P_a(g)=\sum_{m\in\mathbf N^s}c_{a,m}\,g^m,
\]
with finitely many nonzero coefficients $c_{a,m}\in K$, not all zero.
Fix $a$ with $P_a$ not identically zero, and let $m_0$ minimize
$v(c_{a,m})$ over the nonzero coefficients. Then all
$c_{a,m}/c_{a,m_0}$ lie in $R_v$, and the residue of
$P_a(g)/c_{a,m_0}$ is a nontrivial polynomial in the residues
$\bar g_j$ with coefficients in $\kappa(v)$ and one coefficient equal to
$1$. Algebraic independence of the $\bar g_j$ over $\kappa(v)$ makes this
residue nonzero, so
\[
 w\bigl(P_a(g)\bigr)=v(c_{a,m_0})\in\Gamma_v .
\]
Consequently the nonzero summands $f^aP_a(g)$ have values in the cosets
$\langle a,\gamma\rangle+\Gamma_v$, where $\gamma_i=w(f_i)$; rational
independence of the $\gamma_i$ modulo $\Gamma_v$ makes these cosets, hence
the values, pairwise distinct. A sum with a unique term of least value is
nonzero by \cref{lem:strong-triangle-equality}, a contradiction.

The finite generation of $\Gamma_w/\Gamma_v$ and of
$\kappa(w)/\kappa(v)$ in the equality case is proved in
\cite[Chapter~VI, \S10, no.~3]{AC}; we use it as stated and do not
reproduce the proof.
\end{proof}

\begin{corollary}\label{cor:rank-plus-trdeg}
In the setting of \cref{thm:abhyankar-inequality}, one also has
\[
 \operatorname{rank}(v)+\operatorname{trdeg}(v)\leq n .
\]
\end{corollary}
\begin{proof}
Combine \cref{thm:abhyankar-inequality} with
$\operatorname{rank}(v)\leq\operatorname{rat.rk}(v)$ from
\cref{lem:rank-leq-rational-rank}.
\end{proof}

Equality defines an \emph{Abhyankar valuation}.

The refinement in \cref{thm:abhyankar-relative} is not decorative. It
says that an Abhyankar valuation has value group $\mathbf Z^{\,r}$ for
some $r$ and a residue field which is again a function field; both facts
propagate through the gallery of the next section, where the refinement
is used to compute a residue field that resists direct attack. Every
quasimonomial valuation is Abhyankar: by
\cref{lem:monomial-invariants} below, a monomial valuation on an
$n$-dimensional regular model satisfies
$\operatorname{rat.rk}+\operatorname{trdeg}=n$, and all four invariants
depend only on the valuation, not on the model. Over ground fields of
characteristic zero the converse also holds---Abhyankar valuations of a
function field are quasimonomial on suitable log smooth models, by
resolution arguments; see \cite{JM12}---but we never use the converse,
and no monomialization is asserted in positive characteristic.

\begin{example}[Divisorial valuations]\label{ex:divisorial-valuations}
Let $X$ be an $n$-dimensional integral variety and let $E$ be a prime divisor
on a normal birational model $Y\to X$. Then
$\operatorname{ord}_E$ has rational rank $1$ and residue transcendence degree
$n-1$. It is therefore Abhyankar. A positive multiple
$c\operatorname{ord}_E$, $c>0$, is called a divisorial valuation as well.
\end{example}

This is the classical valuation-theoretic convention, which permits an
arbitrary positive real rescaling.  In the trivially valued analytic cone
we use the more restrictive phrase \emph{divisorial point} for a positive
rational multiple of $\operatorname{ord}_E$; the scale is then part of
the analytic point.  The two conventions will always be distinguished.

The assertions in the example are immediate from what we have proved:
since $Y$ is normal, the local ring $\mathcal O_{Y,\eta_E}$ at the generic
point of $E$ is a one-dimensional normal noetherian local domain, hence a
discrete valuation ring, so $\Gamma_{\operatorname{ord}_E}=\mathbf Z$ and
the rational rank is $1$; the residue field is the function field
$k(E)$ of the $(n-1)$-dimensional variety $E$, of transcendence degree
$n-1$ over $k$. Thus divisorial valuations saturate
\cref{thm:abhyankar-inequality} with the maximal possible residue
transcendence degree.  The converse characterization of this extreme case
is a separate theorem of valuation theory and is not needed here.
Divisorial valuations reappear as the
component valuations $b_E^{-1}\operatorname{ord}_E$ of models in
\cref{chap:models}, and, over a trivially valued field, as the divisorial
points whose density and finite-energy properties drive
\cref{part:bj}; see \cref{chap:trivial-analytification} and
\cref{chap:bj-linear-growth}.

\section{A gallery of valuations on
\texorpdfstring{$k(x,y)$}{k(x,y)}}
\label{sec:valuation-gallery}
\index{valuation!gallery of examples}

This section computes every invariant introduced above for four
valuations of the rational function field $F=k(x,y)$, trivial on an
algebraically closed ground field $k$. Here $n=2$, so Abhyankar's
inequality reads
$\operatorname{rat.rk}(v)+\operatorname{trdeg}(v)\leq2$. The specimens
realize the possibilities
\[
 1+1,\qquad 2+0,\qquad 2+0\ \text{(rank two)},\qquad 1+0,
\]
the last one being the strict inequality. We first compute all monomial
valuations at once, by graded reduction.

\begin{lemma}
\label{lem:monomial-invariants}
\index{monomial valuation!invariants}
Let $k$ be a field, $A=k[x_1,\ldots,x_n]$, and
$\alpha\in\mathbf R_{\geq0}^n$. Define
\[
 v_\alpha\Bigl(\sum_\beta c_\beta x^\beta\Bigr)
 =\min_{c_\beta\neq0}\langle\alpha,\beta\rangle ,
\]
and let $L_\alpha=\{\beta\in\mathbf Z^n\mid\langle\alpha,\beta\rangle=0\}$
and $d=\dim_{\mathbf Q}\sum_i\mathbf Q\alpha_i$. Then:
\begin{enumerate}
\item $v_\alpha$ extends to a valuation on $k(x_1,\ldots,x_n)$, trivial
      on $k$, with value group
      $\Gamma_{v_\alpha}=\sum_i\mathbf Z\alpha_i$ and rational rank $d$;
\item the graded reduction $\operatorname{gr}_{v_\alpha}A$ is the
      polynomial ring $k[\bar x_1,\ldots,\bar x_n]$ on the initial forms,
      graded by $\deg\bar x_i=\alpha_i$;
\item the residue field is the fraction field of the group algebra
      $k[L_\alpha]$; it is purely transcendental over $k$, of
      transcendence degree
      \[
       \operatorname{trdeg}(v_\alpha)
       =\operatorname{rank}_{\mathbf Z}L_\alpha=n-d ;
      \]
\item consequently
      $\operatorname{rat.rk}(v_\alpha)+\operatorname{trdeg}(v_\alpha)=n$:
      every monomial valuation, and hence every quasimonomial valuation,
      is Abhyankar.
\end{enumerate}
Moreover the center of $v_\alpha$ on $\mathbf A^n$ is the generic point
of the coordinate subspace $\{x_i=0\mid\alpha_i>0\}$.
\end{lemma}
\begin{proof}
(2) For $\gamma\geq0$, the group
$F^{\geq\gamma}A/F^{>\gamma}A$ has as a $k$-basis the classes of the
monomials $x^\beta$ with $\langle\alpha,\beta\rangle=\gamma$, because the
value of a polynomial is computed monomialwise and distinct monomials are
linearly independent. Multiplication of classes is induced by
multiplication of monomials, which adds weights. Hence
$\operatorname{gr}_{v_\alpha}A$ is the free commutative $k$-algebra on
$\bar x_i=\operatorname{in}(x_i)$ with the stated grading. In particular
it is a domain, so initial forms are multiplicative,
$v_\alpha(fg)=v_\alpha(f)+v_\alpha(g)$, and $v_\alpha$ is a valuation on
$A$ with zero support, extending to $F$ by differences.

(1) The values on $A\setminus\{0\}$ are the weights
$\langle\alpha,\mathbf N^n\rangle$, which generate
$\sum_i\mathbf Z\alpha_i$ as a group; by
\cref{lem:graded-residue-field}(3) this group is $\Gamma_{v_\alpha}$. Its
rational rank is $\dim_{\mathbf Q}\sum_i\mathbf Q\alpha_i=d$ by
definition. Triviality on $k$ is clear.

(3) By \cref{lem:graded-residue-field}(2), $\kappa(v_\alpha)$ is the
degree-zero part of the homogeneous fraction field of
$k[\bar x_1,\ldots,\bar x_n]$. A degree-zero homogeneous fraction is a
quotient $h/g$ of homogeneous elements of equal weight $\gamma$. Fix a
monomial $x^{\beta_0}$ of weight $\gamma$ occurring in $g$; dividing
numerator and denominator by $\bar x^{\beta_0}$ writes $h/g$ as a
quotient of two elements of the group algebra
$k[L_\alpha]
=\bigoplus_{\beta\in L_\alpha}k\,\bar x^{\beta}$, embedded in the
homogeneous fraction field. Conversely every element of
$k[L_\alpha]$ is a degree-zero fraction. Distinct lattice points give
linearly independent monomials, so
$k[L_\alpha]$ is the group algebra of the free abelian group
$L_\alpha\simeq\mathbf Z^{\,n-d}$: choosing a lattice basis identifies it
with a Laurent polynomial ring in $n-d$ variables, whose fraction field
is purely transcendental of degree $n-d$. The rank computation
$\operatorname{rank}L_\alpha=n-d$ holds because
$L_\alpha\otimes\mathbf Q$ is the kernel of the linear map
$\mathbf Q^n\to\mathbf R$, $\beta\mapsto\langle\alpha,\beta\rangle$,
whose image $\sum_i\mathbf Q\alpha_i$ has dimension $d$.

(4) Add the numbers from (1) and (3). For the quasimonomial case, let
$v$ be monomial with weights $\alpha_1,\ldots,\alpha_r>0$ in a regular
system of parameters $z_1,\ldots,z_r$ at a codimension-$r$ point $\xi$ of an
$n$-dimensional regular model $Y$ with function field $F$, and let
$d=\dim_{\mathbf Q}\sum_i\mathbf Q\alpha_i$. The graded computation in
(2) applies verbatim to the monomial expansion from the definition: the
associated graded is a weighted polynomial ring over
$\kappa(\xi)$, a domain, so $v$ is a valuation. We now bound its
invariants on $F$ from below. The values $v(z_i)=\alpha_i$ are realized
in $F$, so $\operatorname{rat.rk}(v|_F)\geq d$. For residues: units of
$\mathcal O_{Y,\xi}$ surject onto $\kappa(\xi)$, whose transcendence
degree over $k$ is $n-r$ because $\xi$ is the generic point of an
$(n-r)$-dimensional subvariety; and for a basis
$\beta^{(1)},\ldots,\beta^{(r-d)}$ of the lattice
$L_\alpha\subset\mathbf Z^r$, the monomial ratios
$z^{\beta^{(j)}}\in F$ are units of $R_v$ whose residues are
algebraically independent over $\kappa(\xi)$, by the graded reduction
argument of (3). Transcendence degrees add in towers, so
$\operatorname{trdeg}(v|_F)\geq(n-r)+(r-d)$. Hence
\[
 \operatorname{rat.rk}(v|_F)+\operatorname{trdeg}(v|_F)
 \geq d+(n-r)+(r-d)=n ,
\]
and \cref{thm:abhyankar-inequality} forces equality. The invariants are
intrinsic to the valuation on the function field, so the conclusion is
independent of the chosen model.

For the center: the prime ideal
$\{f\in A\mid v_\alpha(f)>0\}$ is spanned by the monomials of positive
weight, and a monomial has positive weight exactly when it is divisible
by some $x_i$ with $\alpha_i>0$; so this ideal is
$(x_i\mid\alpha_i>0)$, the generic point of the stated subspace.
\end{proof}

\begin{example}[A divisorial valuation: $\alpha=(1,1)$]
\label{ex:gallery-divisorial}
\index{divisorial valuation!on the plane}
Let $v=v_{(1,1)}$ on $k(x,y)$, so that $v(f)$ is the order of vanishing
of $f$ at the origin: the least total degree in $f$. By
\cref{lem:monomial-invariants}, $\Gamma_v=\mathbf Z$, the graded
reduction is $k[\bar x,\bar y]$ with the total-degree grading,
$L_{(1,1)}=\mathbf Z(1,-1)$, and
\[
 \kappa(v)=k(\bar x/\bar y),
 \qquad
 \operatorname{rat.rk}(v)=1,\quad
 \operatorname{trdeg}(v)=1,\quad
 \operatorname{rank}(v)=1 .
\]
Abhyankar's inequality is an equality, $1+1=2$.

The valuation is divisorial: let $\rho:Y\to\mathbf A^2$ be the blowup of
the origin with exceptional divisor $E$. On the chart with coordinates
$(u,t)$, $\rho(u,t)=(u,ut)$, the divisor $E$ is $\{u=0\}$ and
\[
 \operatorname{ord}_E\bigl(f(u,ut)\bigr)
 =\min_{c_\beta\neq0}\operatorname{ord}_u
 \bigl(u^{\beta_1}(ut)^{\beta_2}\bigr)
 =\min_{c_\beta\neq0}(\beta_1+\beta_2)=v(f),
\]
where the middle minimum is exact because the terms
$c_\beta u^{\beta_1+\beta_2}t^{\beta_2}$ of least $u$-order do not cancel:
their $t$-exponents reproduce the initial form. The residue field
$k(\bar x/\bar y)$ is the function field of
$E\simeq\mathbf P^1$, with $\bar x/\bar y$ the standard coordinate, as
predicted by the divisorial example above. Over a discretely valued
field the analogous computations produce the component valuations of
models in \cref{chap:models}; over the trivially valued $k$, the ray
$c\,v$, $c>0$, is one of the divisorial rays whose density in the
analytification is proved in \cref{chap:trivial-analytification}.
\end{example}

\begin{example}[A quasimonomial valuation with irrational weights:
$\alpha=(1,\sqrt2)$]
\label{ex:gallery-irrational}
\index{quasimonomial valuation!irrational weights}
Let $w=v_{(1,\sqrt2)}$. Since $1$ and $\sqrt2$ are linearly independent
over $\mathbf Q$, every polynomial has a single monomial as initial
form, and \cref{lem:monomial-invariants} gives
\[
 \Gamma_w=\mathbf Z+\sqrt2\,\mathbf Z\simeq\mathbf Z^2,
 \qquad
 L_{(1,\sqrt2)}=0,
 \qquad
 \kappa(w)=k,
\]
so
\[
 \operatorname{rat.rk}(w)=2,\quad
 \operatorname{trdeg}(w)=0,\quad
 \operatorname{rank}(w)=1 ,
\]
the rank being one because $\Gamma_w\subset\mathbf R$ is archimedean
(\cref{prop:rank-one-real-embedding}). Abhyankar's inequality is again an
equality, $2+0=2$, now with the value group carrying all the information.

This valuation is not divisorial, for the intrinsic reason that its
rational rank is $2$: no blowup can convert it into an
$\operatorname{ord}_E$. Geometrically, $w$ lies in the one-dimensional
normalized slice $v(x)=1$ of the two-dimensional monomial cone attached to
the SNC pair $(\mathbf A^2,\{xy=0\})$: the monomial valuations
$v_{(1,s)}$, $s>0$, trace that slice between the two divisorial rays of the coordinate
axes, rational $s$ giving points that become divisorial after rescaling
and toroidal blowup, irrational $s$ giving points that never do. This is
the local picture of the skeleta of \cref{chap:skeleta} and of the
quasimonomial simplices of \cref{chap:trivial-analytification}, and it is
the two-dimensional analogue of the type II/type III dichotomy of
\cref{ex:graded-gauss}.
\end{example}

\begin{example}[A rank-two lexicographic valuation]
\label{ex:gallery-lex}
\index{valuation!lexicographic}\index{rank!two example}
Order $\mathbf Z^2$ lexicographically:
$(i,j)<(i',j')$ if $i<i'$, or $i=i'$ and $j<j'$. Define
\[
 v_{\mathrm{lex}}\Bigl(\sum c_{ij}x^iy^j\Bigr)
 =\min_{c_{ij}\neq0}{}^{\mathrm{lex}}\ (i,j)\in\mathbf Z^2_{\mathrm{lex}} .
\]
Each pair $(i,j)$ is the weight of exactly one monomial, so the graded
reduction is again $k[\bar x,\bar y]$, now with the
$\mathbf Z^2$-grading; it is a domain, so $v_{\mathrm{lex}}$ is a
valuation on $k(x,y)$, trivial on $k$, with value group
$\mathbf Z^2_{\mathrm{lex}}$.

\emph{Invariants.} The rational rank is $2$. The subgroup
$\Delta=0\oplus\mathbf Z$ is convex: if
$|(i,j)|\leq|(0,m)|$ lexicographically, then $i=0$. The chain
$0\subsetneq\Delta\subsetneq\mathbf Z^2$ shows the rank is $2$; by
\cref{prop:rank-one-real-embedding} the group
$\mathbf Z^2_{\mathrm{lex}}$ is not archimedean---indeed
$n(0,1)<(1,0)$ for every $n$---so $v_{\mathrm{lex}}$ admits no
real-valued normalization. The degree-zero component of the graded
reduction is $k$, so $\kappa(v_{\mathrm{lex}})=k$ and the residue
transcendence degree is $0$. Both inequalities are saturated:
\[
 \operatorname{rat.rk}+\operatorname{trdeg}=2+0=2,
 \qquad
 \operatorname{rank}+\operatorname{trdeg}=2+0=2 ,
\]
the latter being the extreme case of \cref{cor:rank-plus-trdeg}.

\emph{Composite structure.} The valuation factors through the divisor
$\{x=0\}$: for $f\neq0$, write $f=x^af_1$ with $x\nmid f_1$. Then the
first component of $v_{\mathrm{lex}}(f)$ is
$a=\operatorname{ord}_{x=0}(f)$, and the second is the order of vanishing
at $y=0$ of the restriction
$f_1|_{x=0}\in k(y)^\times$. Thus $v_{\mathrm{lex}}$ is the composition
of the divisorial valuation $\operatorname{ord}_{x=0}$ of $k(x,y)$ with
the valuation $\operatorname{ord}_0$ of its residue field $k(y)$. This is
an instance of composite valuations, whose ranks add; see
\cite[Chapter~VI, \S10]{ZariskiSamuelII}.

\emph{Relation to the rank-one world.} The points of our analytic spaces
are real semivaluations, so $v_{\mathrm{lex}}$ is not a point of
$(\mathbf A^2)^{\mathrm{an}}$. It is, however, a limit of monomial
points: for $f=\sum c_{ij}x^iy^j$ and all sufficiently small
$\varepsilon>0$,
\[
 v_{(1,\varepsilon)}(f)=\min_{c_{ij}\neq0}(i+\varepsilon j)
\]
is attained at the lexicographic minimum, so
$v_{(1,\varepsilon)}(f)\to a$, the first component of
$v_{\mathrm{lex}}(f)$, as $\varepsilon\to0^+$, and the rescaled second
coordinate $\varepsilon^{-1}(v_{(1,\varepsilon)}(f)-a)$ stabilizes to the
second component. Higher-rank valuations thus record iterated rates of
degeneration of rank-one families; the systematic home for them is the
Riemann--Zariski space \cite{Tem11}, which we use only as a guiding
picture.
\end{example}

\begin{example}[An infinitely singular valuation]
\label{ex:gallery-infsingular}
\index{valuation!infinitely singular}
Assume $\operatorname{char}k=0$. We construct a rank-one valuation of
$k(x,y)$, trivial on $k$, with
\[
 \operatorname{rat.rk}(v)=1,
 \qquad
 \operatorname{trdeg}(v)=0,
\]
so that Abhyankar's inequality is strict, $1+0<2$; its value group is a
non-finitely-generated subgroup of $\mathbf Q$. In the terminology of the
two-dimensional valuative classification such valuations are called
infinitely singular: they are the valuations of transcendental
``Puiseux series with unbounded denominators.''

\emph{Step 1: a ring of generalized power series.} Let $R$ be the set of
formal sums $a=\sum_{\gamma\in S_a}c_\gamma t^\gamma$ with coefficients
in $k$, exponents in $\mathbf Q_{\geq0}$, and \emph{locally finite}
support: $S_a\cap[0,B]$ is finite for every $B\geq0$. Sums and products are defined by the usual formulae; in a
product, the coefficient of $t^\gamma$ is a finite sum because each
factor has only finitely many exponents below $\gamma$. Define
$\operatorname{ord}_t(a)=\min S_a$ for $a\neq0$. If $a,b\neq0$, the
exponent $\min S_a+\min S_b$ has exactly one representation as a sum of
exponents from $S_a$ and $S_b$, so its coefficient in $ab$ is the product
of the two leading coefficients: thus $R$ is a domain and
$\operatorname{ord}_t$ is a valuation on it, extending to
$\operatorname{Frac}(R)$. Compare the examples of generalized power
series valuations in \cite[Chapter~VI]{ZariskiSamuelII}.

\emph{Step 2: the series and the embedding.} Put
\[
 b_j=j+1-2^{1-j}\qquad(j\geq1),
\]
so $b_1=1$, the sequence is strictly increasing and unbounded, and $b_j$
has exact denominator $2^{j-1}$: indeed
$b_j=\bigl((j+1)2^{j-1}-1\bigr)/2^{j-1}$ with odd numerator for
$j\geq2$. Let
\[
 \hat y=\sum_{j\geq1}t^{b_j}\in R,
 \qquad
 s_\ell=\sum_{j\leq\ell}t^{b_j},
 \qquad
 r_\ell=\hat y-s_\ell ,
\]
so $s_\ell\in k[t^{1/m_\ell}]$ with $m_\ell=2^{\ell-1}$, and
$\operatorname{ord}_t r_\ell=b_{\ell+1}$.

We claim $\hat y$ is transcendental over $k(t)$. Suppose
$P(t,Y)\in k[t][Y]$ is nonzero of minimal $Y$-degree $D$ with
$P(t,\hat y)=0$; then $D\geq1$ and, by minimality and
$\operatorname{char}k=0$, the polynomial
$\partial_YP$ does not vanish at $(t,\hat y)$. Set
$c_i=\operatorname{ord}_t\bigl(\partial_Y^iP(t,\hat y)\bigr)$ for
$0\leq i\leq D$, so $c_0=+\infty$,
$\mu:=c_1<+\infty$, and let $C=\min_ic_i$ over the finite values. All
supports involved lie in the group $\mathbf Z[1/2]$, so
$\mu$ has some exact denominator $2^{a}$. Taylor expansion in $R$ gives
the two exact identities
\[
 \begin{aligned}
  0=P(t,\hat y)
   &=P(t,s_\ell)
     +\sum_{i\geq1}\tfrac1{i!}\,\partial_Y^iP(t,s_\ell)\,r_\ell^i,\\
  \partial_Y^iP(t,s_\ell)
   &=\sum_{l\geq0}\tfrac1{l!}\,
     \partial_Y^{i+l}P(t,\hat y)(-r_\ell)^l .
 \end{aligned}
\]
From the second identity,
$\operatorname{ord}_t\partial_Y^iP(t,s_\ell)\geq C$ for all $i,\ell$,
and for $i=1$ the terms with $l\geq1$ have order at least
$C+b_{\ell+1}$; since $b_\ell\to\infty$, for all large $\ell$ we get
$C+b_{\ell+1}>\mu$ and hence
$\operatorname{ord}_t\partial_YP(t,s_\ell)=\mu$ exactly. In the first
identity, the term $i=1$ then has order exactly $\mu+b_{\ell+1}$, while
each term with $i\geq2$ has order at least $C+2b_{\ell+1}>\mu+b_{\ell+1}$
for large $\ell$. By the strong triangle equality,
\[
 \operatorname{ord}_tP(t,s_\ell)=\mu+b_{\ell+1}
 \qquad\text{for all large }\ell .
\]
But $P(t,s_\ell)$ lies in $k[t^{1/m_\ell}]$, so its order lies in
$2^{-(\ell-1)}\mathbf Z$, whereas $\mu+b_{\ell+1}$ has exact denominator
$2^{\ell}$ once $2^\ell$ exceeds $2^{a}$: a contradiction. Hence no such
$P$ exists, and
\[
 x\mapsto t,\qquad y\mapsto\hat y
\]
embeds $k[x,y]$ into $R$. Let $v$ be the restriction of
$\operatorname{ord}_t$ along this embedding: a valuation on $k(x,y)$,
trivial on $k$, with $v(x)=1$ and $v(y)=b_1=1$.

\emph{Step 3: the value group is not finitely generated.} Fix
$\ell\geq1$, write $q(u)=\sum_{j\leq\ell}u^{e_j}\in k[u]$ with
$e_j=b_jm_\ell=(j+1)2^{\ell-1}-2^{\ell-j}\in\mathbf N$, so that
$s_\ell=q(t^{1/m_\ell})$, and set
\[
 M_\ell(x,Y)
 :=\prod_{\zeta^{m_\ell}=1}\bigl(Y-q(\zeta u)\bigr)\Big|_{u^{m_\ell}=x} .
\]
The coefficients of the product are polynomials in $u$ invariant under
$u\mapsto\zeta'u$ for every $m_\ell$-th root of unity $\zeta'$, hence
polynomials in $u^{m_\ell}$; so $M_\ell\in k[x,Y]$, and
$M_\ell(x,y)\in k[x,y]$. Evaluating in $R$,
\[
 v\bigl(M_\ell(x,y)\bigr)
 =\sum_{\zeta^{m_\ell}=1}
 \operatorname{ord}_t\bigl(\hat y-q(\zeta t^{1/m_\ell})\bigr).
\]
For $\zeta=1$ the order is $\operatorname{ord}_tr_\ell=b_{\ell+1}$. For
$\zeta$ a primitive $2^i$-th root of unity, $1\leq i\leq\ell-1$, note
$e_j\equiv-2^{\ell-j}\pmod{2^i}$, which is nonzero modulo $2^i$ exactly
when $j>\ell-i$; hence
\[
 \hat y-q(\zeta t^{1/m_\ell})
 =\sum_{j\leq\ell}(1-\zeta^{e_j})t^{b_j}+r_\ell
\]
has order $b_{\ell-i+1}$, the smallest $b_j$ whose coefficient survives.
There are $2^{i-1}$ such roots, so
\[
 v\bigl(M_\ell(x,y)\bigr)
 =b_{\ell+1}+\sum_{i=1}^{\ell-1}2^{i-1}b_{\ell+1-i}.
\]
The term $b_{\ell+1}$ has exact denominator $2^{\ell}$, while each
summand $2^{i-1}b_{\ell+1-i}$ lies in $2^{-(\ell-1)}\mathbf Z$; hence
$v(M_\ell(x,y))$ has exact denominator $2^{\ell}$. As $\ell$ grows, the
value group $\Gamma_v\subset\mathbf Z[1/2]$ therefore contains elements
of unbounded denominator and is not finitely generated. For instance
$M_1=y-x$ with $v(y-x)=b_2=\tfrac52$, and
$M_2=(y-x)^2-x^5$ with value $b_3+b_2=\tfrac{15}4+\tfrac52=\tfrac{25}4$.

\emph{Step 4: invariants.} Since
$\Gamma_v\subset\mathbf Q$, the group is archimedean, so
$\operatorname{rank}(v)=1$ and $\operatorname{rat.rk}(v)=1$. If the
residue transcendence degree were $1$, Abhyankar's inequality would be an
equality and \cref{thm:abhyankar-relative} would force $\Gamma_v$ to be
finitely generated, contradicting Step 3. Hence
$\operatorname{trdeg}(v)=0$ and
\[
 \operatorname{rat.rk}(v)+\operatorname{trdeg}(v)=1<2 :
\]
the valuation is not Abhyankar, hence not quasimonomial on any model, by
\cref{lem:monomial-invariants}. It is nevertheless a perfectly good
rank-one point of the analytification. Points of this kind are unavoidable in
the valuative spaces of this book, and the global theory of
\cref{part:bj} is designed to handle them through limits along nets
rather than through models.
\end{example}

The following table summarizes the gallery; $\Gamma$ denotes the value
group and the last column records
$\operatorname{rat.rk}+\operatorname{trdeg}$ against $n=2$.
\[
\begin{array}{l|c|c|c|c|c}
 &\Gamma&\operatorname{rank}&\operatorname{rat.rk}
 &\operatorname{trdeg}&\text{Abhyankar}\\
\hline
 v_{(1,1)}=\operatorname{ord}_E
 &\mathbf Z&1&1&1&1+1=2\\
 v_{(1,\sqrt2)}
 &\mathbf Z+\sqrt2\mathbf Z&1&2&0&2+0=2\\
 v_{\mathrm{lex}}
 &\mathbf Z^2_{\mathrm{lex}}&2&2&0&2+0=2\\
 v_{\text{inf.\ sing.}}
 &\subset\mathbf Q,\ \text{not f.g.}&1&1&0&1+0<2
\end{array}
\]
The first, second, and fourth rows are rank-one real-valued valuations and
hence points of the Berkovich analytification with the displayed
normalizations; only the first is divisorial.  The lexicographic valuation
is not a Berkovich point because it is not real-valued. The fourth row is
not Abhyankar. The third row is only encoded by an iterated asymptotic and is not a limit
inside the Berkovich space.

\section{Linear growth}

We record a notion needed in the trivially valued theory. Let $X$ be a
projective integral variety, $L$ an ample line bundle, and $v$ a valuation of
$K(X)$ trivial on the ground field. Fix a nonzero rational section $\tau$ of
$L$. For $s\in H^0(X,mL)\setminus\{0\}$ set
\[
                            v_\tau(s)=v(s/\tau^m).
\]

\begin{definition}\label{def:valuation-linear-growth}
The valuation $v$ has \emph{linear growth with respect to $L$} if
\[
  \sup_{m\geq1}\frac{1}{m}
  \sup_{0\ne s\in H^0(X,mL)}v_\tau(s)<+\infty.
\]
\end{definition}

This property is independent of $\tau$: replacing $\tau$ by $g\tau$ subtracts
$m v(g)$ from every $v_\tau(s)$, changing the displayed normalized quantities
by the fixed constant $-v(g)$. It is also independent of replacing $L$ by a
positive tensor power. Divisorial valuations have linear growth; the general
criterion in terms of finite energy is developed in \cref{part:bj}.

We prove the divisorial case, since it anchors the finite-energy theory of
\cref{chap:bj-linear-growth} and the proof displays the mechanism: an
ample class bounds the multiplicity of any effective divisor along a fixed
prime divisor, linearly in the degree.

\begin{lemma}\label{lem:divisorial-linear-growth}
\index{linear growth!of divisorial valuations}
Let $X$ be a projective integral variety of dimension $n$ over a field
$k$, let $L$ be ample, and let $v=c\operatorname{ord}_E$ be a divisorial
valuation, where $E$ is a prime divisor on a normal projective birational
model $\rho:Y\to X$ and $c>0$. Then $v$ has linear growth with respect to
$L$.
\end{lemma}
\begin{proof}
Fix a nonzero rational section $\tau$ of $L$ and an ample line bundle $A$
on $Y$. For $0\neq s\in H^0(X,mL)$, the divisor
$D_s:=\operatorname{div}(s)\in|mL|$ is effective and
\[
 v_\tau(s)=v(s/\tau^m)
 =c\operatorname{ord}_E\bigl(\rho^*D_s\bigr)
 -mc\operatorname{ord}_E\bigl(\rho^*\operatorname{div}(\tau)\bigr).
\]
The second term is $m$ times a constant independent of $s$. For the
first, write $b=\operatorname{ord}_E(\rho^*D_s)\geq0$. The cycle
$\rho^*D_s-bE$ is effective, and intersecting with $A^{n-1}$ gives
\[
 b\,(A^{n-1}\cdot E)\leq(A^{n-1}\cdot\rho^*D_s)
 =m\,(A^{n-1}\cdot\rho^*L),
\]
where the equality uses that intersection numbers against line bundles
depend only on the linear equivalence class of $D_s$. The degree
$(A^{n-1}\cdot E)$ is strictly positive because $E$ is an integral
$(n-1)$-dimensional projective scheme and $A$ is ample. Hence
\[
 b\leq m\,\frac{(A^{n-1}\cdot\rho^*L)}{(A^{n-1}\cdot E)},
\]
and $\sup_{s}v_\tau(s)\leq Cm$ for a constant $C$ independent of $m$.
The intersection theory used is that of proper schemes over a field
\cite[Chapter~2]{Ful}.
\end{proof}

\begin{example}[Monomial valuations have linear growth]
\label{ex:monomial-linear-growth}
Let $X=\mathbf P^n_k$ with homogeneous coordinates $X_0,\ldots,X_n$, let
$L=\mathcal O(1)$, and let $v_\alpha$ be the monomial valuation of
\cref{lem:monomial-invariants} in the affine coordinates $x_i=X_i/X_0$,
with weights $\alpha\in\mathbf R_{\geq0}^n$. Take
$\tau=X_0$. A section $s\in H^0(X,m L)$ is a degree-$m$ form, and
$s/X_0^m$ is a polynomial of degree at most $m$ in $x_1,\ldots,x_n$, so
\[
 v_\alpha(s/X_0^m)
 =\min_{c_\beta\neq0}\langle\alpha,\beta\rangle
 \leq m\max_i\alpha_i ,
\]
since each exponent vector $\beta$ appearing has $|\beta|\leq m$.
Hence $v_\alpha$ has linear growth, with constant at most
$\max_i\alpha_i$. The trivial valuation ($\alpha=0$) is the degenerate
case: all $v_\tau(s)$ vanish. By contrast, let a semivaluation have nonzero
support and be centered in a chart on which a chosen local frame $\tau$ is
nonvanishing.  If the local coefficient of a section $s$ belongs to that
support, then $v(s/\tau^m)=+\infty$, so no finite linear-growth bound can
hold.  This is the only support-bearing formulation used here; no quotient
is evaluated when its denominator lies in the support. The finite-energy interpretation of the linear-growth locus
occupies \cref{chap:bj-linear-growth}.
\end{example}

\section{A dictionary for later use}

The same object will often be written in two coordinate systems:
\[
\begin{array}{c|c}
\text{multiplicative notation}&\text{additive notation}\\ \hline
|f(x)|&v_x(f)=-\log|f(x)|\\
\max(|f|,|g|)&\min(v(f),v(g))\\
|T(x)|=\rme^{-r}&v_x(T)=r\\
\text{radius }\rho&\text{value }-\log\rho.
\end{array}
\]
Every later formula involving slopes, model functions, or tropicalization is
checked against this table.

\section*{Source notes}

The valuation-theoretic foundations are classical; we use
\cite{ZariskiSamuelII} and the Stacks Project for scheme-theoretic valuative
criteria. Berkovich points and completed residue fields follow
\cite[Chapter~1]{Berk12}. Graded reduction in this form is due to Temkin;
see \cite[\S1]{Temkin2004}.
Abhyankar valuations and their monomialization are
subtle in positive characteristic; no local uniformization theorem is being
asserted here. Quasimonomial means \emph{monomial on an already available
regular model}; it does not assert monomialization of every Abhyankar
valuation in every characteristic.
